% ID: SEL-D0032-OA-E02
% ESTADO: validada. S033, cotejada con examen y criterios.
\textbf{a)} Cada tramo es continuo en su dominio.
En $6$, el primero vale $-90+120=30$ y el límite por la derecha
es $(540-240)/10=30$. Por tanto la ocupación es continua
en todo $[0,\infty)$, incluida la continuidad lateral en $0$.

\textbf{b)} Dos años son 24 meses. En el segundo tramo,
$f(24)=1920/28=480/7$, aproximadamente $\boxed{68.57\%}$.

\textbf{c)} En el primer tramo, igualar a 40 da
$t^2-8t+16=0$, es decir, $(t-4)^2=0$.
En el segundo, $90t-240=40(t+4)$ implica $t=8$.
Ambas raíces pertenecen a su tramo. Se alcanza el $40\%$
a los $\boxed{4\text{ y }8\text{ meses}}$.

\textbf{d)} El primer tramo se escribe $40-\frac52(t-4)^2$
y no supera 40. El segundo se escribe $90-600/(t+4)$
y es siempre inferior a 90. A largo plazo tiende al $90\%$,
así que el complejo nunca llega al $100\%$.
